% Source: https://artofproblemsolving.com/community/c3h1772214_cheater_shopkeeper
% Stable identity: aops-topic-1772214-q01
Suppose the true price of the goods is $1$ dollar per kilogram. When buying, he pays for $1$ kg but receives $1.1$ kg, so $1.1$ kg cost him $1$ dollar.

When selling, each time he hands over $1$ kg he charges for $1.1$ kg, i.e.\ $1.1$ dollars. So the $1.1$ kg he bought for $1$ dollar sell for $1.1\cdot1.1=1.21$ dollars. His profit is $0.21$ dollars on a cost of $1$ dollar, that is, $\boxed{21\%}$.

(If the cheating when selling is read instead as handing over $10\%$ less than the buyer pays for, $0.9$ kg for the price of $1$ kg, the profit is $\dfrac{1.1}{0.9}-1=\dfrac{2}{9}=22\dfrac{2}{9}\%$. The statement above fixes the first reading.)
